\documentclass[a4paper,10pt]{article}

\usepackage{parskip}
\usepackage{enumitem}
\usepackage[a4paper,margin=0.8in]{geometry}
\usepackage{titlesec}
\usepackage{hyperref}
\thispagestyle{empty}

\usepackage{carlito}
\renewcommand{\familydefault}{\sfdefault}

\setlength{\parindent}{0pt}

\titlespacing{\section}{0pt}{0.6em}{0.4em}
\titleformat{\section}{\bfseries\fontsize{12}{14}\selectfont}{}{0em}{}
[\titlerule]

\setlist[itemize]{
  leftmargin=12pt,
  itemsep=-5pt,
  topsep=2pt
}

\begin{document}

{\fontsize{14}{16}\selectfont \textbf{Rohan Batra}} \\
\vspace{-10pt}
\fontsize{11}{13}\selectfont

\href{mailto:contact@rohanbatra.in}{contact@rohanbatra.in} \textbar{}
\href{tel:+917417879921}{+91-74178-79921} \textbar{}
\href{https://rohanbatra.in}{rohanbatra.in} \textbar{}
\href{https://linkedin.com/in/rohanbatrain}{linkedin.com/in/rohanbatrain} \textbar{} \\
\href{https://github.com/rohanbatrain}{github.com/rohanbatrain}

\vspace{-2pt}

\section*{Professional Summary}

\fontsize{10}{12}\selectfont
Computer Science undergraduate (UPES, '27) focused on reliability and platform engineering --- building and operating production systems on Linux and cloud infrastructure, and root-causing incidents end-to-end. I design and operate UPES-ECS, an always-on emergency-communication system on enterprise networking (active/standby HA, health checks, load testing, runbooks). HACKIITK 2026 finalist at IIT Kanpur (top $\sim$0.4\% of 1,300+ teams). Comfortable across Linux, AWS, Docker, CI/CD, and security audits (OWASP, SAST/DAST) throughout the SDLC.

\section*{Education}

\fontsize{10}{12}\selectfont
B.Tech in Computer Science Engineering, UPES Dehradun \hfill Aug 2027 \\
CGPA: 8/10

\section*{Technical Skills}

\fontsize{10}{12}\selectfont
\textbf{Reliability \& Systems:} Linux, systemd, High availability / failover, Health checks, Incident response, Load testing, Runbooks, Proxmox \\
\textbf{Cloud, DevOps \& Automation:} AWS (EC2, ECS, S3), Google Cloud, Docker, GitHub Actions, CI/CD, System Automation \\
\textbf{Networking:} Juniper (SRX, EX, Mist), TCP/IP, DNS, VLAN, QoS \\
\textbf{Software Engineering:} Python, C/C++, JavaScript, SQL, FastAPI, REST APIs, Git \\
\textbf{Security \& Accessibility:} OAuth2, RBAC, OWASP Top 10, SAST/DAST, WCAG 2.0 \\
\textbf{AI \& Tooling:} LangChain, Pydantic, Ollama, Redis \\
\textbf{Web \& Mobile:} Hugo, Flutter, Android SDK, HTML, CSS

\section*{Experience}

\textbf{Techymonks, Bangalore (Remote)} \hfill Nov 2025 – Present \\
\textit{Software Engineer Intern} \\
\vspace{-17pt}
\begin{itemize}
    \item Building Python/FastAPI backend services deployed on AWS (EC2, ECS, S3)
    \item Conducting application security audits combining manual code review, automated scanning (SAST/DAST), and vulnerability reproduction against OWASP Top 10
    \item Integrating automated security checks into CI/CD pipelines as part of the continuous integration workflow
    \item Contributing to WCAG 2.0 accessibility on frontend integrations and backend-supported features
\end{itemize}

\textbf{Google Developer Groups On Campus (GDGOC), UPES Dehradun} \hfill Sep 2025 – Present \\
\textit{Technical Head} \\
\vspace{-17pt}
\begin{itemize}
    \item Led and delivered technical workshops for 100+ participants
    \item Managed cross-functional teams executing production-grade initiatives
    \item Mentored students in cloud, system design, and full-stack development
\end{itemize}

\textbf{LinuxWale, Dehradun} \hfill Sep 2023 – Apr 2025 \\
\textit{Linux Support \& Community Volunteer} \\
\vspace{-17pt}
\begin{itemize}
    \item Onboarded 10+ students to Linux: distribution selection, dual-boot setup, and driver/hardware troubleshooting across varied laptops
    \item Wrote shell scripts to automate post-install configuration and routine setup tasks
    \item Mentored peers transitioning from Windows to Linux as a daily-driver operating system
\end{itemize}

\textbf{Cybersecurity Awareness Sessions, Indian Armed Forces Cantonment Welfare Programs} \hfill Aug 2024 – Sep 2024 \\
\textit{Cybersecurity Awareness Trainer} \\
\vspace{-17pt}
\begin{itemize}
    \item Delivered cybersecurity and social media safety awareness sessions within Indian Armed Forces cantonment premises across three official welfare programs
    \item Engaged with families of Indian Armed Forces personnel and senior officers through practical sessions on digital safety, privacy risks, online fraud prevention, and responsible social media usage
    \item Adapted technical cybersecurity concepts for diverse audiences, including spouses, children, and non-technical participants
    \item Recognized with three Certificates of Appreciation from Golden Key Eagles Women Empowerment Center, HQ Uttarakhand Sub Area, and Golden Key Tuskers
\end{itemize}
\textbf{Uttarakhand Blind Sports Association, Vikasnagar (On-site)} \hfill Jun 2024 – Jul 2024 \\
\textit{Student Volunteer} \\
\vspace{-17pt}
\begin{itemize}
    \item Supported sports programs for visually impaired participants
    \item Coordinated logistics and on-ground event execution
    \item Contributed to accessibility-focused initiatives
\end{itemize}

\section*{Projects}

\textbf{UPES-ECS — Campus Emergency Communication System} \hfill 2026 \\
Showcased to Team HPE \textbar{} \href{https://rohanbatrain.github.io/UPES-ECS/}{rohanbatrain.github.io/UPES-ECS} \\
Tech Stack: Linux, Asterisk/PJSIP (VoIP/SIP), Python/FastAPI, Docker, HPE Juniper (SRX/EX/Mist) \\
\vspace{-17pt}
\begin{itemize}
    \item Designed and operate a resilient, offline-first campus emergency communication system --- one number (111) reaches a trained response team with no internet, cloud, or cellular dependency
    \item Runs on an enterprise HPE + Juniper networking backbone (SRX firewall, EX2300 PoE switches, Mist Wi-Fi 6) with VLAN/QoS prioritization and high-availability failover
    \item Built a Python/FastAPI operations console and telemetry API with an automated call flow, offline voice coaching, and 44-language localized prompts generated offline
    \item Presented live to Team HPE during their campus visit; the architecture extends beyond campus to disaster-response and critical-communication zones
\end{itemize}

\textbf{HACKIITK Finale — Gap Hunter (Team Hunterz)} \hfill 2026 \\
Organized by C3iHub, IIT Kanpur \\
Tech Stack: Python, LangChain, Ollama, Pydantic, Electron, GitHub Actions \\
\vspace{-17pt}
\begin{itemize}
    \item Built an offline AI-powered policy gap analyzer mapping security policies against NIST CSF 2.0 and the CIS MS-ISAC framework, producing evidence-grounded gap reports, revised policy drafts, and remediation roadmaps
    \item Designed a multi-stage pipeline (PDF/DOCX ingestion $\rightarrow$ hybrid retrieval combining vector + sparse + reranking $\rightarrow$ gap analysis $\rightarrow$ policy revision) orchestrated via LangChain with Pydantic-validated structured outputs to ground LLM reasoning and reduce hallucination
    \item Shipped cross-platform Electron desktop app with local Ollama runtime integration; automated Linux/macOS/Windows releases via GitHub Actions and PyInstaller
    \item Maintained 702-test suite (97.7\% pass rate) spanning unit, integration, property-based, and adversarial/chaos tests
    \item Placed top $\sim$0.4\% of 1,300+ teams and 9,000+ participants in the Solutions Track; selected as one of 24 finalist teams invited to IIT Kanpur campus
\end{itemize}

\textbf{Second Brain Database} \hfill Nov 2024 \\
Tech Stack: Python, FastAPI, Redis, OAuth2, RBAC, Docker \\
\vspace{-17pt}
\begin{itemize}
    \item Built secure REST APIs using OAuth2 authentication and role-based access control (RBAC)
    \item Designed optimized database schemas and integrated Redis for caching
    \item Containerized services using Docker for consistent CI/CD deployments
\end{itemize}

\textbf{n8n Nodes — Second Brain Database} \hfill Jan 2025 \\
Tech Stack: JavaScript, n8n, Workflow Automation \\
\vspace{-17pt}
\begin{itemize}
    \item Developed custom n8n nodes exposing Second Brain Database APIs as reusable workflow steps
    \item Enabled end-to-end automation pipelines integrating backend services with orchestration workflows
\end{itemize}

\textbf{Portfolio \& Resume CI/CD Automation} \hfill Jan 2025 \\
Tech Stack: Hugo, Flutter, MarkLang, GitHub Actions \\
\vspace{-17pt}
\begin{itemize}
    \item Built multi-source deployment pipeline merging a Hugo blog and a Flutter portfolio into a unified \texttt{gh-pages} deployment
    \item Integrated MarkLang as Hugo's content translation layer for blog posts, bridging a custom language with a static site generator
    \item Automated build, merge, and deploy on commit; eliminated manual release steps for both portfolio and resume
\end{itemize}

\textbf{MarkLang (Translation Engine)} \hfill Jan 2025 \\
\vspace{-17pt}
\begin{itemize}
    \item Built an Ollama-powered translation engine for automated multilingual workflows
    \item Automated the full translation pipeline using GitHub Actions for autonomous execution
\end{itemize}

\textbf{Proxmox Auto-Install Assistant} \hfill Jan 2025 \\
\vspace{-17pt}
\begin{itemize}
    \item Automated installation of Proxmox, reducing setup time by $\sim$80\%
    \item Simplified complex system configuration workflows through scripted provisioning
\end{itemize}

\textbf{Karmstrot Builds (Android \& C++)} \hfill Aug 2024 \\
Tech Stack: C++, Shell, Android SDK, KernelSU \\
\vspace{-17pt}
\begin{itemize}
    \item Built automation scripts for Android KernelSU integration across devices
    \item Modified low-level system configurations and applied SDLC practices for versioning and testing
\end{itemize}

\section*{Certifications \& Accomplishments}

\fontsize{10}{12}\selectfont
\begin{itemize}
    \item Top 15\% — TryHackMe cybersecurity learning platform
    \item Certificate of Appreciation — Golden Key Eagles Women Empowerment Center (Sep 2024)
    \item Certificate of Appreciation — HQ Uttarakhand Sub Area (Aug 2024)
    \item Certificate of Appreciation — Golden Key Tuskers, AWWA Week 2024 (Aug 2024)
    \item Semifinalist — Eureka! Junior 2022--23, E-Cell IIT Bombay (Issued Nov 2022; Credential ID: ecell2021)
    \item Semifinalist — Eureka! Junior 2021--22, E-Cell IIT Bombay (Issued Dec 2021; Credential ID: ecell2021)
\end{itemize}

\end{document}
